% Short communication for Mathematical Programming Computation
% (svjour3, smallextended; at most 10 pages in the submitted version).
\RequirePackage{fix-cm}
\documentclass[smallextended]{svjour3}
\smartqed

\usepackage{amsmath,amssymb}
\usepackage{mathtools}
\usepackage{graphicx}
\usepackage{booktabs}
\usepackage{xcolor}
\usepackage{listings}
\usepackage{url}

\lstset{language=Python,basicstyle=\ttfamily\footnotesize,columns=fullflexible,
  keepspaces=true,commentstyle=\color{gray},frame=none,xleftmargin=1em,
  aboveskip=4pt,belowskip=2pt}

\journalname{Mathematical Programming Computation}

\newcommand{\R}{\mathbb{R}}
\newcommand{\N}{\mathbb{N}}
\newcommand{\GL}{\mathrm{GL}}
\newcommand{\Sym}{\mathbb{S}}
\newcommand{\SOS}{\Sigma}
\newcommand{\cB}{\mathcal{B}}
\newcommand{\cI}{\mathcal{I}}
\newcommand{\cQ}{\mathcal{Q}}
\newcommand{\lmin}{\lambda_{\min}}
\newcommand{\lmax}{\lambda_{\max}}
\newcommand{\sz}[2]{s(#1,#2)}
\newcommand{\code}[1]{\texttt{#1}}
\DeclareMathOperator{\Tr}{Tr}

\begin{document}

\title{Combining state lifting with matrix and correlative sparsity in
moment-SOS relaxations of polynomial matrix inequalities}

\titlerunning{State lifting and sparsity in relaxations of PMIs}

\author{Jakub Mare\v{c}ek}

\institute{J. Mare\v{c}ek \at
  Faculty of Electrical Engineering, Czech Technical University in Prague,
  Prague, Czech Republic\\
  \email{jakub.marecek@fel.cvut.cz}}

\date{Received: date / Accepted: date}

\maketitle

\begin{abstract}
The matrix moment-sum-of-squares (SOS) hierarchy of Henrion and Lasserre,
Scherer and Hol, and Kojima solves polynomial matrix inequalities (PMIs) and
polynomial matrix optimization problems globally, but its semidefinite
programs grow as $m\binom{n+d}{d}$ for an $m\times m$ PMI in $n$ variables
at relaxation order~$d$. Chordal sparsity of the matrix, correlative sparsity
of the variables and term sparsity reduce this size, but when every entry of
the matrix depends on every variable, the blocks still grow rapidly with~$n$. For scalar
problems, Balada Gaggioli, Henrion and Korda showed that composition and
low-rank structure can then be exploited by lifting intermediate states to
new variables. We extend their state-lifting hierarchies to matrix-valued
objectives and combine them with the chordal decomposition of the matrix.
Since correlative sparsity alone need not converge for matrix-valued
objectives, we prove that the push-forward variant does. We explain the
combination by a simple observation: the techniques restrict different
factors of the index set of the Gram matrix of a matrix-SOS certificate, so
that their savings multiply. All of this is implemented in \code{pysparsepmi}, an open-source Python package
with a CVXPY-style modelling interface, which also recovers and certifies
minimizers from the pseudo-moments. On a PMI of degree~$80$ in $40$
variables, where the dense relaxation would need a Gram matrix of order
$10^{23}$, the lifted relaxations are tight, and the push-forward one is
solved in 74~seconds.
\keywords{Polynomial matrix inequalities \and Moment-SOS hierarchy \and
Chordal sparsity \and Correlative sparsity \and Term sparsity \and State
lifting \and Semidefinite programming}
\subclass{90C22 \and 90C23 \and 90C26 \and 65K05}
\end{abstract}

% ----------------------------------------------------------------------
\section{Introduction}
\label{sec:intro}

A polynomial matrix inequality (PMI) asks that a symmetric matrix $F(x)$,
whose entries are polynomials in $x\in\R^n$, be positive semidefinite. PMIs
arise in fixed-order and static output feedback control
\cite{HenrionSebekKucera2003,HenrionLasserre2006,HenrionLofbergKocvaraStingl2005},
in robust stability analysis \cite{HenrionLasserre2012}, and in structural
and topology optimization \cite{Kocvara2021,HandaTyburecKocvara2025,HandaEtAl2025}.
Their natural optimization
form is polynomial matrix optimization (PMO),
\begin{equation}
\label{eq:pmo}
\lambda^\star \;=\; \inf_{x\in K}\; \lmin\bigl(F(x)\bigr),\qquad
K=\{x\in\R^n:\ g_j(x)\ge 0,\ j=1,\dots,p\},
\end{equation}
since $F\succeq 0$ on $K$ if and only if $\lambda^\star\ge 0$.

Two families of methods are available. Nonlinear semidefinite programming
solvers, notably PENNON and PENBMI of Ko\v{c}vara and Stingl
\cite{KocvaraStingl2003,HenrionLofbergKocvaraStingl2005} and PENLAB
\cite{PENLAB}, handle large BMIs and PMIs but find local solutions.
The matrix moment-SOS hierarchy of Kojima \cite{Kojima2003}, Henrion and
Lasserre \cite{HenrionLasserre2006} and Scherer and Hol \cite{SchererHol2006}
gives a monotone sequence of lower bounds on~$\lambda^\star$, converging
under an Archimedean assumption, together with certificates; it is
available in GloptiPoly~3 \cite{GloptiPoly3}, YALMIP \cite{YALMIP2004} and
TSSOS \cite{TSSOS2021}, and it converges finitely under classical
optimality conditions \cite{HuangNie2025}. Its price is size: the order-$d$ relaxation of
\eqref{eq:pmo} has a positive semidefinite (PSD) block of order
$m\,\sz{n}{d}$, $\sz{n}{d}=\binom{n+d}{d}$, with $d\ge\lceil\deg F/2\rceil$.

Four kinds of structure have been exploited to reduce this size.
(i) \emph{Matrix sparsity}: if the sparsity graph of $F$ is chordal, the
certificate splits into blocks indexed by its maximal cliques, either via
PSD matrix completion with auxiliary variables
\cite{KimKojimaMevissenYamashita2011} or directly at the level of SOS
matrices \cite{ZhengFantuzzi2023}.
(ii) \emph{Correlative sparsity} restricts each multiplier to a clique of
variables \cite{Waki2006,Lasserre2006}, also for polynomial SDPs and matrix
constraints \cite{KojimaMuramatsu2009,NieQuTangZhang2026}.
(iii) \emph{Term sparsity} block-diagonalizes Gram matrices along the
supports of the data \cite{TSSOS2021,CSTSSOS2022}; it has
been used for PMIs in structural design \cite{HandaTyburecKocvara2025}.
Miller, Wang and Guo \cite{MillerWangGuo2026} extended term, correlative and
matrix sparsity to PMO, combined them in TSSOS, and showed that correlative
sparsity, unlike in the scalar case, need not converge for matrix-valued
objectives.
(iv) \emph{Composition and low-rank structure}: Balada Gaggioli, Henrion
and Korda \cite{BGHK2025,BGHK2026} lift the intermediate states of a
polynomial evaluated through a chain of low-dimensional maps---dynamical
systems, Markov chains, tensor trains, low-rank (CP) polynomials---to new
variables. This creates correlative sparsity where the original problem has
none, with clique sizes independent of~$n$; so far only for scalar
objectives.

For PMIs whose entries depend on all variables, such as products of
parameter-dependent matrices, (i)--(iii) leave PSD blocks that grow rapidly
with~$n$. Our contributions are:
\begin{enumerate}
\item State lifting for PMIs (Sect.~\ref{sec:lifting}): matrix-valued
terminal maps in both state-lifting hierarchies of \cite{BGHK2026}, combined
with the chordal decomposition of the terminal matrix, and a proof that the
push-forward variant converges, in contrast with correlative sparsity for
matrix objectives \cite{MillerWangGuo2026}.
\item A unifying view (Sect.~\ref{sec:combine}): each technique selects a
family of subsets of the Gram index set $\N^n_d\times[m]$. Matrix sparsity
acts on the factor $[m]$, correlative sparsity on $\N^n_d$, so their
combination is a product of families and the savings multiply; lifting is a
reformulation that produces correlative sparsity.
\item The open-source package \code{pysparsepmi} (Sect.~\ref{sec:software})
with a CVXPY-style interface and moment-side recovery and certification of
minimizers.
\item Experiments (Sect.~\ref{sec:experiments}) showing PSD blocks
independent of the number of variables and of the matrix size, against the
dense, matrix-sparse and term-sparse relaxations.
\end{enumerate}

% ----------------------------------------------------------------------
\section{Background}
\label{sec:background}

Let $\SOS^m[x]$ denote the $m\times m$ SOS matrices,
$S(x)=(z_d(x)\otimes I_m)^\top Q\,(z_d(x)\otimes I_m)$ with $Q\succeq 0$,
where $z_d(x)$ is the vector of the $\sz{n}{d}$ monomials $x^\alpha$,
$\alpha\in\N^n_d$, of degree at most~$d$. The order-$d$ Scherer--Hol
relaxation of \eqref{eq:pmo} is
\begin{equation}
\label{eq:sh}
\lambda_d=\max\Bigl\{t:\ F-tI_m=S_0+\textstyle\sum_{j=1}^p g_jS_j,\
S_j\in\SOS^m[x],\ \deg(g_jS_j)\le 2d\Bigr\}.
\end{equation}
Equalities $h=0$ enter with free polynomial-matrix multipliers and matrix
constraints $G(x)\succeq0$ with the Scherer--Hol pairing $\langle S,G\rangle$
\cite{SchererHol2006}. Then $\lambda_d\le\lambda_{d+1}\le\lambda^\star$,
and $\lambda_d\to\lambda^\star$ when the quadratic module of the $g_j$ is
Archimedean \cite{SchererHol2006,Kojima2003,HenrionLasserre2006}. The dual
of \eqref{eq:sh} is a moment relaxation with matrix-valued pseudo-moments
$Y_\gamma\in\Sym^m$ \cite{HenrionLasserre2006}.

\emph{Matrix sparsity.} Let $C_1,\dots,C_q$ be the maximal cliques of a
chordal extension of the graph on $[m]$ with an edge $\{i,k\}$ whenever
$F_{ik}\not\equiv0$, and $E_C\in\{0,1\}^{|C|\times m}$ select the rows in
$C$. Replacing every $S_j$ in \eqref{eq:sh} by
$\sum_k E_{C_k}^\top S_{j,k}E_{C_k}$, $S_{j,k}\in\SOS^{|C_k|}[x]$, gives
blocks of order $|C_k|\,\sz{n}{d}$. This is the polynomial analogue of the
decomposition of PSD matrices with chordal sparsity
\cite{VandenbergheAndersen2015};
Zheng and Fantuzzi \cite{ZhengFantuzzi2023} proved the corresponding
sparse Positivstellens\"atze, in particular a sparse matrix version of
Putinar's theorem for PMIs strictly feasible on Archimedean sets.

\emph{Correlative sparsity.} Let $V_1,\dots,V_r$ be the maximal cliques of
a chordal extension of the graph on $[n]$ with an edge $\{a,b\}$ whenever
$x_a,x_b$ occur in a common monomial of $F$ or in a common constraint.
Then $S_0=\sum_aS_{0,a}$ with $S_{0,a}\in\SOS^m[x_{V_a}]$ and every $g_j$
receives a multiplier in the variables of a clique that contains its own
\cite{Waki2006,Lasserre2006,KojimaMuramatsu2009}. For scalar objectives this
converges under the running intersection property \cite{Lasserre2006}; for
matrix-valued $F$ it may not \cite[Ex.~5.1]{MillerWangGuo2026}.

\emph{Term sparsity.} TSSOS \cite{TSSOS2021} builds a
graph on the Gram index set whose edges are the pairs of monomials whose
product lies in the (extended) support, alternates support extension and
chordal extension, and keeps one PSD block per clique; the bounds increase
with the number $s$ of steps. For PMO, the nodes are pairs $(\alpha,i)$
\cite{MillerWangGuo2026}.

\emph{State lifting.} Suppose that the objective is computed by a chain
\begin{equation}
\label{eq:chain}
s_1=F_1(x_1),\quad s_i=F_i(s_{i-1},x_i)\ (1<i<N),\quad
p(x)=F_N(s_{N-1},x_N),
\end{equation}
with states $s_i\in\R^{r_i}$ and local variables $x_i\in K_i\subset\R^{k_i}$.
Even when $p$ is dense, adding the $s_i$ as variables with equality
constraints $s_i=F_i(s_{i-1},x_i)$ makes the problem correlatively sparse,
with cliques $\{s_{i-1},x_i,s_i\}$ of $r_{i-1}+k_i+r_i$ variables if the
chordal extension follows the chain (SL-chord; for rank-$r$ CP polynomials,
$r+2$ variables: LRPOP) \cite{BGHK2025,BGHK2026}. Alternatively, polynomial
potentials $V_i(s_i)$ are pushed along the chain with one certificate per
stage (SL-push) \cite{BGHK2026}. Both hierarchies were introduced for scalar
objectives~$p$.

% ----------------------------------------------------------------------
\section{Combining structures}
\label{sec:combine}

\subsection{Families of Gram index sets}

Write $S_0$ of \eqref{eq:sh} through its Gram matrix $Q$ indexed by
$\cI=\N^n_d\times[m]$:
$S_0(x)=\sum_{(\alpha,i),(\beta,k)\in\cI}Q_{(\alpha,i),(\beta,k)}\,
x^{\alpha+\beta}e_ie_k^\top$. A family $\cB$ of subsets of $\cI$ defines
the cone of SOS matrices with Gram matrices $Q=\sum_{B\in\cB}P_B^\top
Q_BP_B$, $Q_B\succeq0$, where $P_B$ selects the rows in~$B$. Every such cone
lies in $\SOS^m[x]$, so every family gives a valid bound; families differ
in the sizes $|B|$, which govern the cost, and in their convergence. With $\N^V_d$ the exponents of degree at
most $d$ supported on $V\subseteq[n]$:
\begin{align*}
\cB_{\rm dense}&=\{\N^n_d\times[m]\}, &
\cB_{\rm mat}&=\{\N^n_d\times C_k\}_k, &
\cB_{\rm cs}&=\{\N^{V_a}_d\times[m]\}_a,
\end{align*}
and $\cB_{\rm ts}$ is the set of cliques of the term-sparsity graph
on~$\cI$. Matrix and correlative sparsity are \emph{rectangular}: they
restrict different factors of $\cI$. Their meet
$\cB_{\rm mat}\wedge\cB_{\rm cs}=\{B\cap B'\}=\{\N^{V_a}_d\times C_k\}_{k,a}$
therefore has blocks of order $|C_k|\,\sz{|V_a|}{d}$: the savings of the two
techniques multiply, $m\,\sz{n}{d}\to\omega_{\rm mat}\,\sz{\omega_{\rm
cs}}{d}$, where $\omega$ denotes the largest clique. Localizing multipliers
are treated in the same way. Term sparsity is not rectangular; it refines
the blocks of either family. Table~\ref{tab:sizes} summarizes the block
orders.

\begin{table}[t]
\caption{Largest PSD block of the certificate at order~$d$. $\omega_{\rm
mat}$, $\omega_{\rm cs}$: largest matrix and variable cliques; for lifted
chains, $\rho=\max_i(r_{i-1}+k_i+r_i)$.}
\label{tab:sizes}
\centering\small
\begin{tabular}{@{}lll@{}}
\toprule
Relaxation & Largest block & Reference\\
\midrule
dense (Scherer--Hol) & $m\,\sz{n}{d}$ & \cite{SchererHol2006,HenrionLasserre2006}\\
matrix sparsity & $\omega_{\rm mat}\,\sz{n}{d}$ & \cite{ZhengFantuzzi2023}\\
correlative sparsity & $m\,\sz{\omega_{\rm cs}}{d}$ & \cite{KojimaMuramatsu2009}\\
term sparsity & data dependent, $\le m\,\sz{n}{d}$ & \cite{MillerWangGuo2026}\\
matrix $\wedge$ correlative & $\omega_{\rm mat}\,\sz{\omega_{\rm cs}}{d}$ & \cite{MillerWangGuo2026}, this paper\\
SL-chord $\wedge$ matrix & $\omega_{\rm mat}\,\sz{\rho}{d}$ & this paper\\
SL-push $\wedge$ matrix & $\omega_{\rm mat}\,\sz{r_{N-1}+k_N}{d}$ & this paper\\
\bottomrule
\end{tabular}
\end{table}

The combination is not free: the meet is a restriction of both families, and
already $\cB_{\rm cs}$ alone may lose convergence for matrix-valued
objectives \cite[Ex.~5.1]{MillerWangGuo2026}. We therefore pair every
relaxation with an a~posteriori certificate (Sect.~\ref{sec:moments}), and
prove convergence for the lifted push-forward hierarchy below, which avoids
correlative sparsity in the matrix certificate. Miller, Wang and Guo combine
matrix sparsity of the \emph{constraints} with correlative sparsity; here
the decomposition of the objective $F$ is combined with correlative
sparsity and lifting.

\subsection{State lifting for PMIs}
\label{sec:lifting}

Let the last map in \eqref{eq:chain} be matrix-valued,
$F(x)=F_N(s_{N-1},x_N)\in\Sym^m$, and let $\lambda^\star=\min_{x\in
K_1\times\dots\times K_N}\lmin(F(x))$. Let $B_i\subset\R^{r_i}$ be boxes
with $F_i(B_{i-1}\times K_i)\subseteq B_i$, computed by interval
arithmetic. Lifting gives the equivalent PMO
\begin{equation}
\label{eq:lifted}
\min\ \lmin\bigl(F_N(s_{N-1},x_N)\bigr)\ \ \text{s.t.}\ \
s_i=F_i(s_{i-1},x_i),\ s_i\in B_i,\ x_i\in K_i,
\end{equation}
whose correlative sparsity graph is a chain.

\emph{SL-chord for PMIs.} We apply \eqref{eq:sh} to \eqref{eq:lifted} with
the family $\cB_{\rm mat}\wedge\cB_{\rm cs}$, where the matrix cliques $C_k$
come from the pattern of $F_N$ and the variable cliques from the
stage-by-stage elimination order of \cite{BGHK2025,BGHK2026}. Each
constraint gets one multiplier per matrix clique, on the smallest variable
clique containing it. All blocks have order at most $\omega_{\rm
mat}\,\sz{\rho}{d}$, independently of $N$ and~$m$.

\emph{SL-push for PMIs.} Let $\cQ_i$ be the truncated quadratic module in
$\R[s_{i-1},x_i]$ generated by the constraints of $K_i$ and $B_{i-1}$, and
$\cQ^m_N$ its matrix version, with the matrix-sparse decomposition. Find
$t$ and potentials $V_i\in\R[s_i]$, $\deg V_i\le\lfloor 2d/\deg F_i
\rfloor$, such that
\begin{equation}
\label{eq:push}
V_1(F_1)-t\in\cQ_1,\quad V_i(F_i)-V_{i-1}\in\cQ_i\ (1<i<N),\quad
F_N-V_{N-1}I_m\in\cQ^m_N .
\end{equation}

\begin{proposition}
\label{prop:valid}
Every feasible $t$ of the SL-chord or SL-push relaxation satisfies
$t\le\lambda^\star$. If the quadratic modules of the $K_i$ are Archimedean
(e.g., boxes), the optimal values $t_d$ of \eqref{eq:push} converge to
$\lambda^\star$ as $d\to\infty$.
\end{proposition}

\begin{proof}
Validity of SL-chord holds because every term of the certificate is PSD at
feasible points of \eqref{eq:lifted}. For SL-push, let $x$ be feasible with
states $s_i\in B_i$; then
$\lmin(F(x))\ge V_{N-1}(s_{N-1})\ge\dots\ge V_1(s_1)\ge t$.
For convergence, fix $\delta>0$ and let
$U_{N-1}(s)=\min_{x\in K_N}\lmin(F_N(s,x))$, continuous on $B_{N-1}$.
By the Stone--Weierstrass theorem there is a polynomial $V_{N-1}$ with
$U_{N-1}-2\delta\le V_{N-1}\le U_{N-1}-\delta$ on $B_{N-1}$, so
$F_N-V_{N-1}I\succeq\delta I$ on $B_{N-1}\times K_N$. Backwards for
$i=N-1,\dots,2$, let $U_{i-1}(s)=\min_{x\in K_i}V_i(F_i(s,x))$ and choose
$V_{i-1}$ with $U_{i-1}-2\delta\le V_{i-1}\le U_{i-1}-\delta$ on $B_{i-1}$;
then $V_i(F_i)-V_{i-1}\ge\delta$ on $B_{i-1}\times K_i$. Finally
$t=\min_{x\in K_1}V_1(F_1(x))-\delta$. All conditions in \eqref{eq:push}
hold strictly on compact sets with Archimedean modules, so they have
certificates of some finite degree by Putinar's theorem \cite{Putinar1993}
and, for the last stage, the matrix Positivstellensatz
\cite{SchererHol2006} in its chordal form \cite{ZhengFantuzzi2023}. With
the value functions $W_{N-1}=U_{N-1}$,
$W_{i-1}(s)=\min_{x\in K_i}W_i(F_i(s,x))$, induction gives
$U_i\ge W_i-2(N-1-i)\delta$ on $B_i$, hence $t\ge\lambda^\star-(2N-1)\delta$.
\qed
\end{proof}

SL-chord uses correlative sparsity in a matrix certificate, which may fail
to converge \cite[Ex.~5.1]{MillerWangGuo2026}; we rely on the a~posteriori
certificate of Sect.~\ref{sec:moments} instead. SL-push needs no
correlative sparsity: its stages are coupled through scalar potentials.

\subsection{Exploiting symmetry under PMI}
In this subsection we will consider the following optimization problem: \begin{equation}\label{eq:pop_pmi_cons}
  \min_{x\in\R^{n}} f(x) \text{ s.t. } F(x)\succeq 0.
\end{equation}
where $f(x)$ is a polynomial whereas $F(x)\in\R[x]^{m\times m}$ is a symmetric $m\times m$ polynomial matrix. The feasible set consists of points $x\in\R^{n}$ for which the polynomial matrix $F(x)$ is PSD. We assume that the problem \eqref{eq:pop_pmi_cons} enjoys some symmetric property. Let's present formally the necessary definitions.

Let $\tau\in S_n$ be a permutation, we let $\tau$ act on the space of polynomial as following: \begin{equation}
  \tau\cdot f(x_1,\dots,x_n) \coloneqq f(x_{\tau^{-1}(1)},\dots,x_{\tau^{-1}(n)}).
\end{equation}
Under this definition, consider a monomial $x^{\alpha}$ where $\alpha=(\alpha_1,\dots,\alpha_n)$, there is \begin{equation}
  \tau \cdot x^{\alpha} = x_1^{\alpha_{\tau(1)}}\dots x_n^{\alpha_{\tau(n)}}.
\end{equation}
Moreover, consider two polynomials $f,g$, there is $\tau\cdot(fg) = \tau \cdot f \tau \cdot g$. Let $G\leq S_n$ be a subgroup, we say that a polynomial $f$ is $G$-invariant if $\forall \tau\in G, \tau\cdot f = f$. Let $A\in\R[x]^{m\times p}$ be a polynomial matrix, with abuse of notation, $\tau \cdot A(x)$ stands for the entry-wise action of $\tau$ to $A(x)$.

Given \eqref{eq:pop_pmi_cons}, from now on we assume that there exists a subgroup $G\leq S_n$ such that \begin{itemize}
  \item the polynomial $f$ is $G$-invariant,
  \item there exists a representation $\phi: G\rightarrow \GL(m)$ such that $\forall \tau\in G, \tau\cdot F(x) = \phi(\tau) F(x) \phi(\tau)^{\mathsf{T}}$.
\end{itemize}
Under this assumption, the objective function $f$ and the feasible set $K:=\{x\in\R^{n}:F(x)\succeq 0\}$ are invariant under coordinate permutations of $G$. Note that the polynomial matrix $F$ need not be invariant under $G$. The existence of representation $\phi$ ensures that $K$ is invariant (geometrically under coordinate permutation) without requiring the (algebraic) invariance of $F(x)$ under $G$.

Let $d \in \N$. Define the vector of monomials $v_d(x)^{\mathsf{T}} = (1,x_1,\dots,x_n,\dots,x_n^{d})\in\R[x]^{s(d)}$ where $s(d)$ is the number of monomials upto degree $d$. There exists a representation $\Psi_d: G \rightarrow \GL(s(d))$ where \begin{equation}
  \forall \tau\in G, \tau\cdot v_d(x) = \Psi_d(\tau) v_d(x),
\end{equation}
similarly $\tau\cdot v_d(x)$ should be understood as entry-wise action of $\tau$ on $v_d(x)$.

Let $R(x)\in\R[x]^{m\times m}$ be a degree $d$ polynomial matrix, there always exists (non-unique) symmetric Gram matrix $S\in \R^{ms(d)\times ms(d)}$ where $R(x) = (v_d(x)\otimes I_m)^{\mathsf{T}} S (v_d(x)\otimes I_m)$ We call $R(x)$ sum of squares if there exists a PSD Gram matrix $S$. Recall the usual Frobenius inner product $\langle A,B\rangle=\Tr(AB)$ for symmetric matrices, let $R(x)$ be a polynomial matrix, we compute \begin{align}
    &\tau\cdot\langle R(x),F(x)\rangle=\tau\cdot\langle (v_d(x)\otimes I_m)^{\mathsf{T}} S (v_d(x)\otimes I_m),F(x)\rangle\\
    =&\tau\cdot\langle S, (v_d(x)\otimes I_m)F(x)(v_d(x)\otimes I_m)^{\mathsf{T}}\rangle\label{eq:RF_1}\\
    =&\left\langle S, \tau\cdot\left(v_d(x)v_d(x)^{\mathsf{T}} \otimes F(x) \right) \right\rangle\label{eq:RF_2} \\
    =& \left\langle S, \tau \cdot v_d(x)\tau \cdot v_d(x)^{\mathsf{T}} \otimes \tau \cdot F(x) \right\rangle\label{eq:RF_3}\\
    =&\left\langle S, \left(\Psi_d(\tau) v_d(x)v_d(x)^{\mathsf{T}}\Psi_d(\tau)^{\mathsf{T}}\right) \otimes \left(\phi(\tau)F(x)\phi(\tau)^{\mathsf{T}}\right) \right\rangle \\
    =& \left\langle S, \left(\Psi_d(\tau)\otimes \phi(\tau)\right)(v_d(x)v_d(x)^{\mathsf{T}} \otimes F(x))(\Psi_d(\tau)\otimes \phi(\tau))^{\mathsf{T}} \right\rangle \label{eq:RF_4}\\
    =& \left\langle\left(\Psi_d(\tau)\otimes \phi(\tau)\right)^{\mathsf{T}} S \left(\Psi_d(\tau)\otimes \phi(\tau)\right), v_d(x)v_d(x)^{\mathsf{T}} \otimes F(x) \right\rangle
\end{align}
Some comments are in order. Equation \eqref{eq:RF_1} used cyclic property of trace. \eqref{eq:RF_2} used linearity of $\tau\cdot$. \eqref{eq:RF_3} used the fact that the action distributes over product and definition of Kronecker product. \eqref{eq:RF_2} and \eqref{eq:RF_4} used the property $(AC)\otimes (BD)=(A\otimes B)(C\otimes D)$. Now we can give the following reduction theorem. Let's \begin{theorem}
  Suppose there is an equation of the form $f-t = \sigma_{0}+\langle R(x),F(x)\rangle$ where $t\in\R$, $\sigma_0,R$ are SOS and $\deg(\sigma_0),\deg(R(x)F(x))\leq 2d$, then it's sufficient to take symmetric $\sigma_0$ and $R(x)$.
\end{theorem}
\begin{proof}
  Denote $\delta=\lceil\deg(F)/2\rceil$ and denote by $\Phi_{d}=\Psi_d\otimes \phi$ the tensor product representation. Write $\sigma_0 = v_d(x)^{\mathsf{T}}S_0v_d(x)$ and $R(x) = (v_{d-\delta}(x)\otimes I_m)^{\mathsf{T}} S (v_{d-\delta}(x)\otimes I_m)$ Take the group average $\frac{1}{|G|}\sum_{\tau\in G} \tau\cdot $ of both sides and use symmetry of $f,F$ : \begin{equation}
    \begin{split}
    f(x)-t =& \left\langle\frac{1}{|G|}\sum_{\tau\in G}\Psi_{d}(\tau)^{\mathsf{T}}S_0\Psi_{d}(\tau), v_d(x)v_d(x)^{\mathsf{T}}\right\rangle\\
    &+\left\langle \frac{1}{|G|}\sum_{\tau\in G}\Phi_{d-\delta}(\tau)^{\mathsf{T}}S\Phi_{d-\delta}(\tau), v_d(x)v_d(x)^{\mathsf{T}} \otimes F(x)\right\rangle
    \end{split}
  \end{equation}
  Denote $S_{0,G}=\frac{1}{|G|}\sum_{\tau\in G}\Psi_{d}(\tau)^{\mathsf{T}}S_0\Psi_{d}(\tau)$,$S_{G}=\frac{1}{|G|}\sum_{\tau\in G}\Phi_{d-\delta}(\tau)^{\mathsf{T}}S\Phi_{d-\delta}(\tau)$. We note that if an equation of the form $f-t=\sigma_0+\langle R(x),F(x)\rangle$, then it's sufficient to take $\sigma_0,R$ whose Gram matrix is the group average $S_{0,G},S_{G}$. Indeed, since $S_{0,G},S_G$ are group averages, they satisfy: \begin{equation}
    \Psi_d(\tau)^{\mathsf{T}}S_{0,G}\Psi_d(\tau) = S_{0,G}, \Phi_{d-\delta}(\tau)^{\mathsf{T}}S_{G}\Phi_{d-\delta}(\tau)=S_{G},\forall \in G.
  \end{equation}
\end{proof}

\subsection{Recovering and certifying minimizers}
\label{sec:moments}

The dual multipliers of the coefficient-matching constraints of
\eqref{eq:sh} for the diagonal entries give, after trace normalization,
scalar pseudo-moments $y_\gamma=\operatorname{tr}Y_\gamma/
\operatorname{tr}Y_0$. If the optimal moments are those of a point mass at
$x^\star$ with weight $uu^\top$, $u$ a unit eigenvector of $F(x^\star)$ for
$\lmin$, then $y_{e_j}=x^\star_j$; several minimizers are extracted from a
flat moment matrix as in \cite{HenrionLasserre2005}. In lifted relaxations, the moments of the
original variables are read off the lifted (SL-chord) or per-stage (SL-push)
measures. Whatever the relaxation, a feasible $\hat x$ with
$\lmin(F(\hat x))-\lambda_d\le\varepsilon$ certifies that $\hat x$ is
$\varepsilon$-optimal and that the bound is within $\varepsilon$ of
$\lambda^\star$, without any convergence theory.

% ----------------------------------------------------------------------
\section{The software}
\label{sec:software}

\code{pysparsepmi} is written in pure Python on top of CVXPY
\cite{CVXPY2016}, so any CVXPY conic solver with PSD cones can be
used (SCS \cite{SCS2016}, Clarabel \cite{Clarabel2026}, MOSEK). Its
modelling layer mirrors CVXPY's LMI interface: \code{polyvar} creates
variables, \code{PolyMatrix} symmetric polynomial matrices whose
coefficients may be numbers or affine CVXPY expressions, and \code{P >> 0}
returns a matrix-SOS constraint with the chordal decomposition (for scalars,
an SOS constraint with correlative sparsity, as YALMIP's \code{sos.csp}
\cite{Lofberg2009}); \code{Problem(objective, constraints).solve()}
mirrors \code{cvxpy.Problem}. High-level functions cover~\eqref{eq:pmo}:
\code{pmi\_optimize} with scalar and matrix inequalities (\code{ineqs}),
equalities (\code{eqs}), matrix sparsity (default), correlative sparsity
(\code{cs=True}) and term sparsity (\code{ts="block"} or \code{"MD"}, as in
TSSOS); \code{sos\_lower\_bound} for scalar problems; and
\code{composition\_lower\_bound(maps, box, method)} for chains
\eqref{eq:chain}, with front-ends for CP and tensor-train polynomials in
TensorLy layouts. Results expose bounds, cliques and block sizes, and
\code{minimizer()}, \code{atoms()}, \code{gap(x)} and
\code{max\_violation(x)} (Sect.~\ref{sec:moments}). The following computes
the worst-case gain of Sect.~\ref{sec:gain}:

\begin{lstlisting}
import pysparsepmi as psp
M = lambda x: [[0.9, 0.5 * x], [0.3 * x, 0.8]]
def last(s, x):    # M(x)' S M(x), S stored as s = (s11, s12, s22)
    S = [[s[0], s[1]], [s[1], s[2]]] if s else [[1, 0], [0, 1]]
    A = M(x)
    return [[sum(A[c][a] * S[c][e] * A[e][b] for c in (0, 1)
                 for e in (0, 1)) for b in (0, 1)] for a in (0, 1)]
def step(s, x):    # lifted state: the three entries of S_i
    R = last(s, x); return [R[0][0], R[0][1], R[1][1]]
res = psp.composition_lower_bound(
    [step] * (n - 1) + [last], box=[(0, 1)] + [1] * (n - 1),
    order=3, method="push", sense="max")
x = res.minimizer()    # maximizer from the pseudo-moments
print(res.value, res.gap(x), res.max_violation(x))
\end{lstlisting}

Chordal extensions use a greedy minimum-degree order unless an elimination
order is prescribed, as for lifted chains, where the greedy order can give
larger cliques. Gram bases are pruned by the diagonal-consistency test of
YALMIP \cite{Lofberg2009}. Lifted states are rescaled by their interval
bounds, $s_i=D_iu_i$; this exact change of variables keeps the data of long
chains of order one and was essential for accuracy. The package, its tests
and the scripts behind all results below are available at
\url{https://github.com/jmarecek/pysparsepmi}.

% ----------------------------------------------------------------------
\section{Computational experiments}
\label{sec:experiments}

All runs used an Apple M2 Pro (32\,GB), Python 3.9, CVXPY 1.7.5 and the
interior-point solver Clarabel 0.11.1 \cite{Clarabel2026} with tolerances
$10^{-10}$, one process per run, a limit of 30 minutes and 16\,GB.
The baselines---the dense relaxation \eqref{eq:sh} (as in GloptiPoly or
TSSOS without sparsity), matrix sparsity \cite{ZhengFantuzzi2023} and the
first step of term sparsity with minimum-degree chordal extension
(\code{TS="MD"} in TSSOS) \cite{MillerWangGuo2026}---run through the same
code and solver, so the comparison isolates the relaxations. Times are in
seconds on a monotonic clock and include the construction of the model in
Python.

\subsection{Worst-case gain of a product of matrices}
\label{sec:gain}

Let $P(x)=M(x_1)\cdots M(x_n)$ with $M(x)=\bigl[\begin{smallmatrix}0.9&0.5x\\
0.3x&0.8\end{smallmatrix}\bigr]$ and $x\in[-1,1]^n$, and consider the
worst-case squared gain $\lambda^\star=\max_x\lmax(P(x)^\top P(x))$ of an
$n$-step linear time-varying system. The $2\times2$ matrix $P^\top P$ has
degree $2n$ and each entry depends on every variable, so neither matrix nor
correlative sparsity applies: the dense relaxation has a block of order
$2\binom{2n}{n}$. Since $P$ is affine in each $x_i$,
$\lmax(P^\top P)=\|P\|^2$ is convex in each $x_i$, so $\lambda^\star$ is
attained at a vertex, and enumeration gives $\lambda^\star$ for $n\le 20$.
As $M(-x)=DM(x)D$ with $D=\mathrm{diag}(1,-1)$, the lifted relaxations use
$x_1\in[0,1]$, which leaves $\lambda^\star$ unchanged and makes the maximizer
unique; the baselines keep the full box, whose sign symmetry term sparsity
exploits \cite{MillerWangGuo2026}. The chain lifts the three entries of
$S_i=M(x_i)^\top S_{i-1}M(x_i)$, $S_0=I$, with the matrix-valued terminal map
$S_n$ (listing in Sect.~\ref{sec:software}), so $r_i=3$ and $\rho=7$.

\begin{table}[t]
\caption{Worst-case gain (Sect.~\ref{sec:gain}): $\lambda^\star$ by vertex
enumeration ($^\dagger$: for $n>20$, the value at the recovered maximizer
rounded to the nearest vertex), largest PSD block, total time in seconds and
relative gap $(\text{bound}-\lambda^\star)/\lambda^\star$. The dense block
$2\binom{2n}{n}$ is listed also where the relaxation was not attempted.
Dense and TS bounds are tight whenever solved. $^*$: Clarabel reported
reduced accuracy.}
\label{tab:gain}
\centering\footnotesize\setlength{\tabcolsep}{1.6pt}
\begin{tabular}{@{}rr rr rr rr rr rr@{}}
\toprule
 & & \multicolumn{2}{c}{dense} & \multicolumn{2}{c}{TS, $s=1$} & \multicolumn{2}{c}{SL-push, $d=2$} & \multicolumn{2}{c}{SL-push, $d=3$} & \multicolumn{2}{c}{SL-chord, $d=2$}\\
\cmidrule(lr){3-4}\cmidrule(lr){5-6}\cmidrule(lr){7-8}\cmidrule(lr){9-10}\cmidrule(l){11-12}
$n$ & $\lambda^\star$ & block & time & block & time & time & gap & time & gap & time & gap\\
 & & & & & & \multicolumn{2}{c}{block } & \multicolumn{2}{c}{block } & \multicolumn{2}{c}{block }\\
\midrule
3 & 3.86006 & 40 & 1.1 &  &  &  &  &  &  &  &  \\
4 & 5.96216 & 140 & 193 &  &  &  &  &  &  &  &  \\
\bottomrule
\end{tabular}

\end{table}

Table~\ref{tab:gain} shows the effect of lifting. The dense relaxation and
term sparsity are tight but stop at $n=4$: at $n=5$ the dense block has
order $504$ and both exceed 30 minutes. The lifted relaxations keep blocks
of order $30$, $70$ and $72$ for every $n$, and their time grows about
linearly. SL-push of order~$3$ and SL-chord of order~$2$ are tight for all
$n$; SL-push bounds the gain of a PMI of degree~$80$ in $40$ variables in
74\,s, where the dense block would have order $2\cdot10^{23}$. For $n>20$
tightness is certified by the recovered maximizer $\hat x\approx(1,\dots,1)$
(Sect.~\ref{sec:moments}). SL-push of order~$2$ is the cheapest, but its gap
grows with~$n$.

\subsection{A banded low-rank PMI}
\label{sec:banded}

To exercise all structures at once, let
$F(x)=\sum_{l=1}^{2}A_l\prod_{i=1}^nf_{l,i}(x_i)$ on $[-1,1]^n$ with
tridiagonal PSD $A_l\in\Sym^m$ and univariate quadratics $f_{l,i}$ whose
Bernstein coefficients are $b_0=1$ and $b_1,b_2\in[1+1/n,1+2/n]$, drawn at
random, as in the scalar instances of \cite[Sect.~4.2]{BGHK2025}. Every
$f_{l,i}\ge1$ with equality only at $x_i=-1$, and $F$ is monotone in each
product, so $\lambda^\star=\lmin(A_1+A_2)$ is known. $F$ is tridiagonal, of
degree $2n$ and dense in the variables. Lifting the partial products gives
a chain with $r_i=2$ whose variable cliques have $r+2=4$ variables
\cite{BGHK2025}, and matrix cliques of size~$2$.

%%BANDED-TABLE%%

%%BANDED-DISCUSSION%%

% ----------------------------------------------------------------------
\section{Conclusions}
\label{sec:conclusions}

Matrix, correlative and term sparsity and state lifting act on different
parts of the Gram index set of a matrix-SOS certificate; combining them
multiplies their savings, and lifting turns PMIs that are dense in the
variables but have composition or low-rank structure into problems where
the other techniques apply. With matrix-valued terminal maps, the lifted
hierarchies give tight bounds on PMIs far beyond the reach of the dense,
matrix-sparse and term-sparse relaxations, and the recovered minimizers
certify them. Open questions include conditions under which SL-chord
converges for PMIs, for instance after scalarization with $y^\top F(x)y$
\cite{ZhengFantuzzi2023,MillerWangGuo2026}; combining term sparsity with
correlative sparsity and lifting in \code{pysparsepmi}, as TSSOS does for
PMO \cite{MillerWangGuo2026}; tighter state enclosures than interval
arithmetic; interfaces to SDP solvers that exploit low-rank solutions
\cite{HabibiKocvaraStingl2024}; and a vectorized model construction, which
dominates the run time of small instances.

% Springer requires declarations; fill in before submission, e.g.:
% \section*{Declarations}
% \textbf{Funding} ... \textbf{Conflict of interest} The author declares no
% conflict of interest. \textbf{Code availability} The code and the scripts
% reproducing all results are available at
% \url{https://github.com/jmarecek/pysparsepmi}.

\bibliographystyle{spmpsci}
\bibliography{references}

\end{document}
